% page 629 of the facsimile (image p-628 of the scan)
% block 162.6 x 246.3 mm ; left margin 39.4 mm
\pgc{17.780mm}{0.000mm}{
\l{\cel{53}621}
\l{}
\l{\sou{voltar}. (\sou{netrans.)}\cel{2}I. (\sou{manejo}). Igar (kavalo) jiro-movar}
\l{\cel{3}plurafoye. - II. (\sou{skermo}). Chanjar sua plaso por evitar tu-}
\l{\cel{3}sheso da la espado di sua adverso. - III. (\sou{kartoludo)}.}
\l{\cel{3}Stroko ye qua onu agas omna preni. - DEFIRS.}
\l{}
\l{\sou{volto}. (elektro).Unajo pri la interdifero di la potenciali}
\l{\cel{3}elektrala = 108 unaji elektromagnetala C.G.S. -\cel{1}\sou{Volto inter}-}
\l{\cel{3}\sou{naciona}\cel{1}: tensiono elektrala qua, aplikate a konduktivo di}
\l{\cel{3}qua la rezistiveso egalesas un\cel{1}\sou{ohm}\cel{1}internaciona, genitas fluo}
\l{\cel{3}de un\cel{1}\sou{ampero}\cel{1}internaciona.- DEFIRS.}
\l{}
\l{\sou{voltijar}. (\sou{netrans.}) I. (\sou{kavalk-arto}). Exercar su gimnastike por}
\l{\cel{3}kustumigar su pri saltar sen estribi adsur kavalo qua avancas}
\l{\cel{3}paze, trote, galope. - II. (\sou{akrobat-arto}). Exercar su pri}
\l{\cel{3}dansar sur kordo laxa. - DEFIR.}
\l{}
\l{\sou{voltmetro}. (\sou{tekn.}) Sorto di galvanometro, per qua onu mezuras}
\l{\cel{3}la forco elektromovala o la falo di la potencialo inter du}
\l{\cel{3}punti. - DEFIRS.}
\l{}
\l{\sou{volumino}. (\sou{geom.}) Parto di spaco, quan okupas korpo. - DEFIS.}
\l{}
\l{\sou{voluntar}. (\sou{trans.}) Esar benigna, komplezema, servema (+serviema)}
\l{\cel{3}til konsentar (agor co o to). - EFI.}
\l{}
\l{\sou{volupto}.\cel{2}Plezuro di la sensi, qua igas juar forte. - EFIS.}
\l{}
\l{\sou{voluto}. (\sou{arkit.)}\cel{2}Ornamento spiralforma an la kapitelo di kolono}
\l{\cel{3}(precipue en la klasifiko-grupo \textquotedbl{}ionika\textquotedbl{}), di konsolo, di}
\l{\cel{3}modiliono. - DEFIRS.}
\l{}
\l{\sou{volvar}. (\sou{trans.}) Aplikar rule kozo an-cirkum altra. - L.}
\l{}
\l{\sou{vomar}. (\sou{trans}.) I. Retrojetar konvulso, per sua boko (la materii}
\l{\cel{3}solida o liquida qui esis en la stomako). - II. (\sou{metaf}.)}
\l{\cel{3}Pulsar ek su (ulo mala). EFIS.}
\l{}
\l{\sou{vomero}. (\sou{anat.}) Osto sen-simetraja, dina, quadrilatera, anomala,}
\l{\cel{3}ye qua konsistas la parto supra di la parieto di la naztrui.-}
\l{}
\l{\sou{vomiko}. (\sou{medic.}) Amaso de materio pusoza qua naskas en la pul-}
\l{\cel{3}moni (che la tuberklosiki) e forjetesas per quaza vomo. .}
\l{}
\l{\sou{vomiknuco}. (\sou{bot.}) Bero venenoza ek arboreto de India. L.\cel{1}\sou{strych}-}
\l{\cel{3}\sou{nos nux vomica}. - EFIS.}
\l{}
\l{\sou{vomitorio}. (en\cel{1}\sou{la Roma antiqua}). Larja aperturo en la klozilo}
\l{\cel{3}por lasar pasar la spekteri, en la amfiteatro qua destinesis}
\l{\cel{3}a la ludi publika.(+publicala). - DEFIS}
\l{}
\l{\sou{vorto}. Son-uno artikulata, ek un o plura silabi,uzata por repre-}
\l{\cel{3}zentar ento od eso-maniero. - Skriburo qua reprezentas ta}
\l{\cel{3}son-uno. - DE.}
\l{}
\l{vorticar. (\sou{netrans.}) Forportesar per jirado rapida. - EI.}
}
